\documentclass[12pt,a4paper]{article}

% -------------------------------------------------
% Sprache / Kodierung
% -------------------------------------------------
\usepackage[T1]{fontenc}
\usepackage[utf8]{inputenc}
\usepackage[ngerman]{babel}
\usepackage{lmodern}
\usepackage{microtype}

% -------------------------------------------------
% Mathematik
% -------------------------------------------------
\usepackage{amsmath,amssymb,amsthm,mathtools}

% -------------------------------------------------
% Layout / Grafik / Farben
% -------------------------------------------------
\usepackage[a4paper,margin=2.3cm]{geometry}
\usepackage{booktabs}
\usepackage{xcolor}
\usepackage[hidelinks,unicode]{hyperref}
\pdfstringdefDisableCommands{%
  \def\ü{ue}\def\ö{oe}\def\ä{ae}\def\Ü{Ue}\def\Ö{Oe}\def\Ä{Ae}\def\ss{ss}%
}
\usepackage[most]{tcolorbox}
\usepackage{enumitem}
\usepackage{array}
\usepackage{siunitx}

% -------------------------------------------------
% Farben
% -------------------------------------------------
\definecolor{lückenrot}{RGB}{180,40,40}
\definecolor{bewiesengrün}{RGB}{30,100,50}
\definecolor{warnorange}{RGB}{200,100,10}
\definecolor{hintergrundblau}{RGB}{235,242,250}
\definecolor{hintergrundgelb}{RGB}{255,252,225}
\definecolor{adischviolett}{RGB}{90,20,140}
\definecolor{fehlerrot}{RGB}{200,0,0}
\definecolor{konjviolett}{RGB}{80,0,130}
\definecolor{e8grau}{RGB}{55,55,90}

% -------------------------------------------------
% tcolorbox-Stile
% -------------------------------------------------
\tcbset{
  lücke/.style={colback=red!8, colframe=lückenrot, fonttitle=\bfseries,
    title={Offene Lücke}},
  ergebnis/.style={colback=green!7, colframe=bewiesengrün,
    fonttitle=\bfseries, title={Bewiesenes Ergebnis}},
  warnung/.style={colback=orange!10, colframe=warnorange,
    fonttitle=\bfseries, title={Achtung / Heuristik}},
  adisch/.style={colback=violet!6, colframe=adischviolett,
    fonttitle=\bfseries, title={2-adische Perspektive}},
  kernidee/.style={colback=hintergrundblau, colframe=blue!60!black,
    fonttitle=\bfseries, title={Kernidee}},
  konj/.style={colback=violet!5, colframe=konjviolett, fonttitle=\bfseries,
    title={Konjektur (unbewiesen)}},
  lte/.style={colback=yellow!8, colframe=e8grau, fonttitle=\bfseries,
    title={LTE-Analyse}},
  hauptsatz/.style={colback=green!10, colframe=bewiesengrün!80!black,
    fonttitle=\bfseries\large, title={Stärkster beweisbarer Satz}},
  korrektur/.style={colback=red!4, colframe=fehlerrot!70!black,
    fonttitle=\bfseries, title={Fehlerkorrektur}},
  sensitivitaet/.style={colback=blue!5, colframe=blue!50!black,
    fonttitle=\bfseries, title={Sensitivitätsanalyse}},
}

% -------------------------------------------------
% Theorem-Umgebungen
% -------------------------------------------------
\newtheorem{definition}{Definition}[section]
\newtheorem{proposition}[definition]{Proposition}
\newtheorem{theorem}[definition]{Theorem}
\newtheorem{lemma}[definition]{Lemma}
\newtheorem{korollar}[definition]{Korollar}
\newtheorem{bemerkung}[definition]{Bemerkung}
\newtheorem{hypothese}[definition]{Hypothese}
\newtheorem{satz}[definition]{Satz}

\newenvironment{beweis}{\begin{proof}[Beweis]}{\end{proof}}

% -------------------------------------------------
% Makros
% -------------------------------------------------
\newcommand{\vii}{\nu_2}
\newcommand{\viii}{\nu_3}
\newcommand{\U}{\mathcal{U}}
\newcommand{\N}{\mathbb{N}}
\newcommand{\Z}{\mathbb{Z}}
\newcommand{\R}{\mathbb{R}}
\newcommand{\E}{\mathbb{E}}
\newcommand{\Ztwo}{\mathbb{Z}_2}
\newcommand{\abs}[1]{\lvert#1\rvert}
\newcommand{\atwo}[1]{\lvert#1\rvert_2}
\newcommand{\logF}{\log F(B)}
\newcommand{\muk}{\mu_k}

% -------------------------------------------------
% Dokument
% -------------------------------------------------
\title{%
  Die $m$-Verteilung an C-Ketten-Endpunkten:\\[0.3em]
  Drei Korrekturen und eine neue Forschungsfrage\\[0.5em]
  \large Korrektur-Appendix zu \textit{collatz\_lyapunov\_konditioniert.tex}%
}
\author{Thomas Hoffbauer}
\date{16.\ Juni 2026}

\begin{document}
\maketitle

\begin{abstract}
Dieses Dokument identifiziert und korrigiert drei mathematische Fehler in
\textit{collatz\_lyapunov\_konditioniert.tex} und entwickelt die eigentliche
Forschungsfrage weiter: die bedingte Verteilung
\[
  \mu_k = \Pr\!\bigl(\vii(3^{k+1}m-1) = r \;\big|\; m \text{ stammt aus C-Kette der Länge } k\bigr).
\]

Die drei Fehler sind: (1) $\E[\nu \mid l{=}k] \approx 2$ ist falsch ---
numerisch und strukturell gilt $\E[\nu \mid l{=}k] \approx 3$;
(2) $P(l{=}k) = 2^{-(k+1)}$ ist keine stationäre Bahnverteilung,
sondern eine Startpunktverteilung; (3) ``LTE garantiert $\alpha \geq 2$'' ist
logisch zirkulär und zudem trivial falsch formuliert.

Als neue Ergebnisse werden bewiesen: $3 \mid m$ für alle C-Ketten-Starts
$n_0 \equiv 11 \pmod{12}$; $\nu \geq 2$ gilt trivial für \emph{alle}
C-Ketten-Starts unabhängig von der Residuenklasse; und die $s$-Verteilung
$s = \viii(m)$ ist geometrisch mit $P(s{=}j) = \tfrac{2}{3} \cdot \left(\tfrac{1}{3}\right)^{j-1}$.

Die korrigierte Driftrechnung ergibt $\E[\logF] \approx -0{,}83$, was
\emph{robuster} ist als der fehlerhafte Wert $-0{,}33$.

\medskip
\noindent\textbf{Status:}
\textcolor{bewiesengrün}{§\,2--3 (Divisibilitäts- und Trivialitätsresultate): bewiesen.}
\textcolor{warnorange}{§\,5 (korrigierte Driftrechnung): heuristisch.}
\textcolor{lückenrot}{§\,6 (Ergodizität, stationäre Maßstruktur): offen.}
\end{abstract}

\tableofcontents
\bigskip

\begin{tcolorbox}[warnung, title={Einordnung dieses Dokuments}]
Die Collatz-Vermutung ist unbewiesen.
Dieser Text trennt konsequent zwischen
\textbf{bewiesenen Aussagen} (grün),
\textbf{Heuristiken} (orange)
und \textbf{offenen Lücken} (rot).
Die Fehlerkorrektur betrifft das Vorzeichen und die Größenordnung
des heuristischen Driftterms, nicht die Beweisbarkeit der Collatz-Vermutung.
\end{tcolorbox}

% ====================================================
\section{Korrektur der drei Fehler in \textit{collatz\_lyapunov\_konditioniert.tex}}
\label{sec:korrekturen}
% ====================================================

\subsection{Fehler 1: \texorpdfstring{$\E[\nu \mid l{=}k] \approx 2$}{E[nu|l=k] ≈ 2} ist falsch}

\begin{tcolorbox}[korrektur, title={Fehler 1 --- Falsche Erwartungswert-Schätzung}]
\textbf{Zitat aus dem Dokument} (Bemerkung nach Lemma 2.1):
\begin{quote}
\itshape
In der Grobabschätzung verwenden wir $\E[\nu \mid l{=}k] \approx 2$ für alle $k$.
\end{quote}
und (Tabelle~1, Spalte $\E[\nu \mid l]$):
alle Einträge für $l \geq 1$ lauten $2{,}0$.

\medskip
\textbf{Fehler:} Das Dokument verwechselt die \emph{Minimalschranke}
$\nu \geq 2$ (was korrekt ist, aber trivial --- s.\,u.) mit dem
\emph{Erwartungswert} $\E[\nu \mid l{=}k]$.

\medskip
\textbf{Korrekte Aussage:} Für alle $k \geq 1$ gilt numerisch
(und strukturell, s.\ Abschnitt~\ref{sec:drift_korrektur})
\[
  \E[\nu \mid l{=}k] \;\approx\; 3,
\]
mit der gleichen geometrischen Verteilung auf $\{2, 3, 4, \ldots\}$
wie für $l = 0$.

\medskip
\textbf{Was bewiesen werden müsste:}
Eine exakte Charakterisierung der bedingten Verteilung
$\mu_k = \Pr(\vii(3^{k+1}m - 1) = r \mid l{=}k)$ für die
tatsächliche Menge der $m$-Werte, die aus C-Ketten der Länge $k$ entstehen.
\end{tcolorbox}

\subsection{Fehler 2: \texorpdfstring{$P(l{=}k) = 2^{-(k+1)}$}{P(l=k) = 2^{-(k+1)}} ist keine stationäre Verteilung}

\begin{tcolorbox}[korrektur, title={Fehler 2 --- Fehlende Maßstruktur}]
\textbf{Zitat aus dem Dokument} (Abschnitt 2):
\begin{quote}
\itshape
$P(l{=}k) = 2^{-(k+1)}$ für $k \geq 1$.
\end{quote}
und in Proposition~4.1 (Gesamtdrift):
\begin{quote}
\itshape
$\E[\logF] = P(l{=}0)\cdot\E[\logF \mid l{=}0] + \sum_{k=1}^{\infty} P(l{=}k)\cdot\E[\logF \mid l{=}k]$.
\end{quote}

\medskip
\textbf{Fehler:} $P(l{=}k) = 2^{-(k+1)}$ ist die Verteilung der C-Kettenlänge
für \emph{uniform zufällig} (mod $2^\infty$) gewählte ungerade Startzahlen $n_0$.
Sie ist \emph{keine} stationäre Verteilung des Collatz-Prozesses auf natürlichen
Zahlen.

Konkret: Für die Drift-Rechnung bräuchte man das Maß $\mu$ unter dem die
Blöcke $B_1, B_2, B_3, \ldots$ einer typischen Collatz-Bahn verteilt sind ---
dieses Maß ist unbekannt.

\medskip
\textbf{Korrekte Aussage:} Die Gleichung $P(l \geq k) = 2^{-k}$ gilt als
zahlentheoretischer Fakt (aus $2^{k+1} \mid (n_0 + 1)$) für uniform
gewählte $n_0$. Für den Einsatz in einer Driftrechnung \emph{auf Bahnen}
ist sie eine heuristische Annahme, keine beweisbare Aussage.

\medskip
\textbf{Was bewiesen werden müsste:}
Ergodizität des Collatz-Systems bzgl.\ eines invarianten Maßes, das die
Block-Statistik $P(l{=}k)$ tatsächlich der geometrischen Verteilung annähert.
Dies ist für die Collatz-Abbildung vollständig offen.
\end{tcolorbox}

\subsection{Fehler 3: \texorpdfstring{``LTE garantiert $\alpha \geq 2$''}{``LTE garantiert alpha >= 2''} ist zirkulär und falsch formuliert}

\begin{tcolorbox}[korrektur, title={Fehler 3 --- Missbrauch des LTE}]
\textbf{Zitat aus dem Dokument} (Achtung-Box in Abschnitt~5):
\begin{quote}
\itshape
Der LTE-Mechanismus erzwingt $\nu \geq 2$ an C-Ketten-Endpunkten.
Damit gilt $\alpha \geq 2$, was den kritischen Wert $\alpha^* \approx 1{,}34$
weit unterschreitet.
\end{quote}

\medskip
\textbf{Fehler:} Hier werden drei Dinge vermischt:
\begin{enumerate}[noitemsep]
  \item \textbf{$\nu \geq 2$ ist trivial, nicht LTE:}
        Für $l \geq 1$ und ungerades $m$ ist $3^{l+1}m$ stets ungerade,
        also $3^{l+1}m - 1$ gerade, also $\vii(3^{l+1}m - 1) \geq 1$,
        also $\nu = 1 + \vii(3^{l+1}m-1) \geq 2$.
        Das LTE wird dafür \emph{nicht} benötigt.
  \item \textbf{LTE gibt nur eine \emph{obere} Schranke:}
        Das LTE liefert $\vii(3^{l+1}m - 1) \leq \vii(3m-1) + \vii(l+1)$
        (ein $O(\log l)$-Oberschranke), keine nichttriviale Unterschranke.
  \item \textbf{$\alpha \geq 2$ ist keine Aussage über den Erwartungswert:}
        Aus $\nu \geq 2$ (punktweise) folgt zwar $\alpha = \E[\nu \mid l{\geq}1] \geq 2$,
        aber das ist sinnlos im Kontext der Sensitivitätsanalyse,
        denn der korrekte Wert ist $\alpha \approx 3$, nicht $\alpha \approx 2$.
\end{enumerate}

\medskip
\textbf{Logische Zirkularität:}
Die Sensitivitätsanalyse zeigt $\E[\logF] > 0$ für $\alpha < 1{,}34$.
Daraus folgt $\E[\logF] < 0$ für $\alpha > 1{,}34$.
Da $\alpha \geq 2 > 1{,}34$ trivial gilt, ist die Schlussfolgerung formal korrekt.
Aber die Annahme $\alpha \approx 2$ (statt des korrekten $\alpha \approx 3$)
suggeriert einen deutlich schwächeren negativen Drift ($-0{,}33$ statt $-0{,}83$).

\medskip
\textbf{Korrekte Formulierung:}
$\nu \geq 2$ gilt trivial für alle C-Ketten-Endpunkte ($l \geq 1$),
unabhängig von der Residuenklasse von $n_0$.
Der Erwartungswert erfüllt $\E[\nu \mid l{=}k] \approx 3$ (heuristisch),
mit $P(\nu = j \mid l{=}k) \approx 2^{-(j-1)}$ für $j \geq 2$.
\end{tcolorbox}

% ====================================================
\section{Beweisbarer Satz: \texorpdfstring{$3 \mid m$}{3 | m} für \texorpdfstring{$n_0 \equiv 11 \pmod{12}$}{n0 ≡ 11 (mod 12)}}
\label{sec:drei_teilt_m}
% ====================================================

\subsection{Setup und Notation}

Sei $n_0$ eine ungerade natürliche Zahl mit einer C-Kette der Länge $l \geq 1$.
Dann gilt per Definition:
\[
  n_0 + 1 = 2^{l+1} \cdot m, \quad m \text{ ungerade},
\]
wobei $l = \vii(n_0 + 1) - 1$ und $m = (n_0 + 1)/2^{l+1}$ die ungerade
Komponente von $n_0 + 1$ ist.

\begin{satz}[$3 \mid m$ für $n_0 \equiv 11 \pmod{12}$]
\label{satz:drei_teilt_m}
Sei $n_0 \equiv 11 \pmod{12}$ und sei $l \geq 1$ die C-Kettenlänge von $n_0$.
Dann gilt $3 \mid m$.
\end{satz}

\begin{beweis}
Schreibe $n_0 = 12a + 11$ für ein $a \in \N_0$.
Dann ist
\[
  n_0 + 1 = 12a + 12 = 12(a+1) = 4 \cdot 3 \cdot (a+1) = 2^2 \cdot 3 \cdot (a+1).
\]
Da $l \geq 1$, gilt $2^{l+1} \geq 4$, und es ist
\[
  m = \frac{n_0 + 1}{2^{l+1}} = \frac{12(a+1)}{2^{l+1}} = \frac{3 \cdot 4 \cdot (a+1)}{2^{l+1}}.
\]
Da $m$ nach Voraussetzung ungerade ist (also $2^{l+1}$ die gesamte Zweierpotenz
von $n_0 + 1$ absorbiert), ist der Faktor $\frac{4(a+1)}{2^{l+1}} \in \N$ eine
ungerade ganze Zahl.
Insbesondere ist $m = 3 \cdot \frac{4(a+1)}{2^{l+1}}$, also $3 \mid m$.
\end{beweis}

\begin{bemerkung}
Der Beweis zeigt mehr: Der Faktor $3$ in $m$ kommt ausschließlich von der
Bedingung $12 \mid (n_0 + 1)$, welche äquivalent zu $n_0 \equiv 11 \pmod{12}$
ist (für ungerades $n_0$).
Für $n_0 \not\equiv 11 \pmod{12}$ ist $3 \mid m$ im Allgemeinen falsch
(z.\,B.\ gilt für $n_0 = 7 \equiv 7 \pmod{12}$, $l=2$: $m = 1$, $3 \nmid m$).
\end{bemerkung}

\subsection{Triviales Resultat: \texorpdfstring{$\nu \geq 2$}{nu >= 2} für alle C-Ketten}

\begin{satz}[$\nu \geq 2$ ohne LTE, für beliebige C-Ketten]
\label{satz:nu_geq_2}
Für jedes ungerade $n_0$ mit C-Kettenlänge $l \geq 1$ gilt
$\nu = \vii(3n_l + 1) \geq 2$,
wobei $n_l = 2 \cdot 3^l \cdot m - 1$ der Kettenendpunkt ist.
\end{satz}

\begin{beweis}
Es ist $3n_l + 1 = 3(2 \cdot 3^l \cdot m - 1) + 1 = 2 \cdot 3^{l+1} \cdot m - 2 = 2(3^{l+1}m - 1)$.
Da $3^{l+1}$ und $m$ beide ungerade sind, ist ihr Produkt $3^{l+1}m$ ungerade,
also $3^{l+1}m - 1$ gerade.
Damit gilt $\vii(3^{l+1}m - 1) \geq 1$, und folglich
$\nu = 1 + \vii(3^{l+1}m - 1) \geq 2$.
\end{beweis}

\begin{tcolorbox}[ergebnis, title={Beweisbarer Satz -- Trivialität von $\nu \geq 2$}]
Die Schranke $\nu \geq 2$ an C-Ketten-Endpunkten ($l \geq 1$) ist ein
elementares Paritätsargument: $3^{l+1}m$ ist stets ungerade, also ist
$3^{l+1}m - 1$ stets gerade.
Das LTE wird hierfür \emph{nicht} benötigt.
Für die Frage des Erwartungswerts $\E[\nu \mid l{=}k]$ sagt dieses Argument
nichts aus.
\end{tcolorbox}

% ====================================================
\section{Die \texorpdfstring{$m$}{m}-Verteilung: Bekanntes und Offenes}
\label{sec:m_verteilung}
% ====================================================

\subsection{Schreibweise und Zerlegung}

Sei $n_0 \equiv 11 \pmod{12}$ mit C-Kettenlänge $l \geq 1$,
und sei $m = (n_0 + 1)/2^{l+1}$ (ungerade).
Nach Satz~\ref{satz:drei_teilt_m} gilt $3 \mid m$; wir schreiben
\[
  m = 3^s \cdot m', \qquad s = \viii(m) \geq 1, \quad 3 \nmid m', \quad m' \text{ ungerade}.
\]
Dann ist der relevante Ausdruck am Kettenendpunkt:
\[
  3^{l+1} \cdot m - 1 = 3^{l+1+s} \cdot m' - 1 = 3^N \cdot m' - 1,
  \quad N = l + 1 + s.
\]

\subsection{Was bekannt ist}

\begin{itemize}[leftmargin=2em, itemsep=4pt]
  \item \textbf{$\nu \geq 2$ immer} (Satz~\ref{satz:nu_geq_2}):
        $\vii(3^N \cdot m' - 1) \geq 1$, da $3^N m'$ ungerade.
  \item \textbf{$3 \mid m$ immer} (Satz~\ref{satz:drei_teilt_m}):
        also $s = \viii(m) \geq 1$ für $n_0 \equiv 11 \pmod{12}$.
  \item \textbf{Parität von $\vii(3^N m' - 1)$:}
        $\vii(3^N m' - 1) \geq 2$ genau dann, wenn $3^N m' \equiv 1 \pmod 4$.
        Da $3 \equiv -1 \pmod 4$, gilt $3^N \equiv (-1)^N \pmod 4$, also:
        \[
          \vii(3^N m' - 1) \geq 2 \iff (-1)^N m' \equiv 1 \pmod 4,
        \]
        was von $N \pmod 2$ und $m' \pmod 4$ abhängt.
  \item \textbf{LTE-Oberschranke:}
        Für festes $m' \equiv 1 \pmod 2$, $3 \nmid m'$ gilt
        $\vii(3^N m' - 1) \leq C + \vii(N)$ für eine von $m'$ abhängige
        Konstante $C = \vii(3m'-1)$. Dies ist eine Oberschranke,
        keine Formel für den Erwartungswert.
\end{itemize}

\subsection{Was offen ist}

\begin{tcolorbox}[lücke, title={Offene Frage: Verteilung $\mu_k$}]
Die zentrale unbekannte Größe ist die bedingte Verteilung
\[
  \mu_k(r) \;=\; \Pr\!\Bigl(\vii(3^{k+1}m - 1) = r \;\Big|\;
  m = \tfrac{n_0 + 1}{2^{k+1}},\; n_0 \text{ C-Ketten-Start der Länge } k\Bigr).
\]
Präzise: Was ist die Verteilung der $m$-Werte, die tatsächlich als
Kofaktor $(n_0 + 1)/2^{l+1}$ bei Collatz-Bahnen auftreten?

\medskip
\textbf{Unbekannt:}
\begin{itemize}[noitemsep]
  \item Ist $m$ equidistributed modulo beliebiger Potenzen von $2$?
  \item Ist $\mu_k$ unabhängig von $k$ (wie die Numerik nahelegt)?
  \item Welche höheren Kongruenzbedingungen an $m$ entstehen durch
        die C-Kettenstruktur?
\end{itemize}
\end{tcolorbox}

\subsection{Strukturelle Beobachtung: \texorpdfstring{$s$}{s}-Verteilung}

Numerisch (für $n_0 \leq 10\,000$, $n_0 \equiv 11 \pmod{12}$, $l \geq 1$)
zeigt sich eine auffällig präzise Übereinstimmung:
\[
  \Pr(s = j) \;\approx\; \frac{2}{3} \cdot \left(\frac{1}{3}\right)^{j-1}
  \quad (j \geq 1),
\]
d.\,h.\ $s = \viii(m)$ ist \emph{geometrisch verteilt} mit Parameter $p = 2/3$.

\begin{tcolorbox}[konj]
\textbf{Konjektur ($s$-Verteilung):}
Für uniform zufällig gewählte $n_0 \equiv 11 \pmod{12}$ mit $l \geq 1$
ist $s = \viii\!\bigl((n_0+1)/2^{l+1}\bigr)$ geometrisch verteilt:
$P(s = j) = \tfrac{2}{3} \cdot \left(\tfrac{1}{3}\right)^{j-1}$.

\smallskip
\textbf{Begründung (heuristisch):}
Gegeben $3 \mid m$, gilt $9 \mid m$ mit Wahrscheinlichkeit $1/3$
(wenn $m/3$ zufällig mod $3$ verteilt ist), und so weiter.
Für einen tatsächlichen Beweis wäre Equidistribution von $m$
modulo Potenzen von $3$ nötig.
\end{tcolorbox}

% ====================================================
\section{Numerische Tabellen}
\label{sec:numerik}
% ====================================================

Die folgenden Tabellen basieren auf einer exhaustiven Berechnung
aller $n_0 \equiv 11 \pmod{12}$, $1 \leq n_0 \leq 10\,000$, $l \geq 1$.
Es wurden $833$ solcher Startzahlen gefunden.

\subsection{Erste 25 Fälle}

\begin{table}[h!]
\centering
\caption{C-Ketten-Starts $n_0 \equiv 11 \pmod{12}$: erste 25 Fälle}
\label{tab:erste_faelle}
\renewcommand{\arraystretch}{1.25}
\begin{tabular}{r r r r r r r}
\toprule
$n_0$ & $l$ & $m$ & $s=\viii(m)$ & $m'$ & $N=l{+}1{+}s$ & $\nu$ \\
\midrule
  11 & 1 &  3 & 1 &  1 & 3 &  2 \\
  23 & 2 &  3 & 1 &  1 & 4 &  5 \\
  35 & 1 &  9 & 2 &  1 & 4 &  5 \\
  47 & 3 &  3 & 1 &  1 & 5 &  2 \\
  59 & 1 & 15 & 1 &  5 & 3 &  2 \\
  71 & 2 &  9 & 2 &  1 & 5 &  2 \\
  83 & 1 & 21 & 1 &  7 & 3 &  3 \\
  95 & 4 &  3 & 1 &  1 & 6 &  4 \\
 107 & 1 & 27 & 3 &  1 & 5 &  2 \\
 119 & 2 & 15 & 1 &  5 & 4 &  3 \\
 131 & 1 & 33 & 1 & 11 & 3 &  4 \\
 143 & 3 &  9 & 2 &  1 & 6 &  4 \\
 155 & 1 & 39 & 1 & 13 & 3 &  2 \\
 167 & 2 & 21 & 1 &  7 & 4 &  2 \\
 179 & 1 & 45 & 2 &  5 & 4 &  3 \\
 191 & 5 &  3 & 1 &  1 & 7 &  2 \\
 203 & 1 & 51 & 1 & 17 & 3 &  2 \\
 215 & 2 & 27 & 3 &  1 & 6 &  4 \\
 227 & 1 & 57 & 1 & 19 & 3 & 10 \\
 239 & 3 & 15 & 1 &  5 & 5 &  2 \\
 251 & 1 & 63 & 2 &  7 & 4 &  2 \\
 263 & 2 & 33 & 1 & 11 & 4 &  2 \\
 275 & 1 & 69 & 1 & 23 & 3 &  3 \\
 287 & 4 &  9 & 2 &  1 & 7 &  2 \\
 299 & 1 & 75 & 1 & 25 & 3 &  2 \\
\bottomrule
\end{tabular}
\end{table}

\begin{bemerkung}
Beobachtungen aus Tabelle~\ref{tab:erste_faelle}:
\begin{itemize}[noitemsep]
  \item $3 \mid m$ in \textbf{allen} Fällen --- bestätigt Satz~\ref{satz:drei_teilt_m}.
  \item $\nu \geq 2$ in \textbf{allen} Fällen --- bestätigt Satz~\ref{satz:nu_geq_2}.
  \item $m'$ nimmt viele verschiedene Werte an (keine offensichtliche Kongruenzstruktur).
  \item Für $n_0 = 227$: $\nu = 10$ (ein großer Ausreißer: $3^3 \cdot 19 - 1 = 512 = 2^9$).
\end{itemize}
\end{bemerkung}

\subsection{Häufigkeitstabelle von \texorpdfstring{$\nu$}{nu} nach Kettenlänge}

\begin{table}[h!]
\centering
\caption{Verteilung von $\nu$ nach Kettenlänge $l$ (alle $n_0 \equiv 11 \pmod{12}$, $n_0 \leq 10\,000$)}
\label{tab:nu_verteilung}
\renewcommand{\arraystretch}{1.3}
\begin{tabular}{r r r r r r r r}
\toprule
$l$ & $n$ & $\bar\nu$ & $\nu{=}2$ & $\nu{=}3$ & $\nu{=}4$ & $\nu{=}5$ & $\nu{\geq}6$ \\
\midrule
1 & 417 & 3,002 & 209 & 104 &  52 & 26 & 26 \\
2 & 208 & 3,000 & 104 &  52 &  26 & 13 & 13 \\
3 & 104 & 3,019 &  52 &  26 &  13 &  7 &  6 \\
4 &  52 & 2,942 &  26 &  13 &   7 &  3 &  3 \\
5 &  26 & 3,115 &  13 &   6 &   3 &  2 &  2 \\
6 &  13 & 3,077 &   6 &   3 &   2 &  1 &  1 \\
7 &   7 & 3,143 &   4 &   1 &   1 &  0 &  1 \\
8 &   3 & 3,000 &   1 &   1 &   1 &  0 &  0 \\
9 &   2 & 3,500 &   1 &   0 &   0 &  1 &  0 \\
\bottomrule
\end{tabular}
\end{table}

\begin{tcolorbox}[ergebnis, title={Numerisches Hauptergebnis: $\E[\nu \mid l{=}k] \approx 3$}]
Für alle $k \in \{1, \ldots, 9\}$ und alle $n_0 \equiv 11 \pmod{12}$,
$n_0 \leq 10\,000$:
\[
  \E[\nu \mid l{=}k] \approx 3 \quad \text{(schwankt zwischen } 2{,}94 \text{ und } 3{,}5\text{)}.
\]
Die Verteilung folgt sehr präzise $P(\nu{=}j \mid l{=}k) \approx 2^{-(j-1)}$
für $j \geq 2$ (geometrisch auf $\{2, 3, 4, \ldots\}$).
Insbesondere: $\nu = 1$ \emph{kommt nie vor}.
\end{tcolorbox}

\subsection{Verteilung von \texorpdfstring{$s = \viii(m)$}{s = nu3(m)}}

\begin{table}[h!]
\centering
\caption{Verteilung von $s = \viii(m)$ für alle 833 Fälle}
\label{tab:s_verteilung}
\renewcommand{\arraystretch}{1.3}
\begin{tabular}{r r r r}
\toprule
$s$ & Anzahl & Häufigkeit & $(2/3)\cdot(1/3)^{s-1}$ \\
\midrule
1 & 556 & 0,6675 & 0,6667 \\
2 & 185 & 0,2221 & 0,2222 \\
3 &  62 & 0,0744 & 0,0741 \\
4 &  20 & 0,0240 & 0,0247 \\
5 &   7 & 0,0084 & 0,0082 \\
6 &   2 & 0,0024 & 0,0027 \\
7 &   1 & 0,0012 & 0,0009 \\
\bottomrule
\end{tabular}
\end{table}

Die Übereinstimmung mit der geometrischen Verteilung
$P(s{=}j) = \tfrac{2}{3}(1/3)^{j-1}$ ist auf vier Nachkommastellen genau.

% ====================================================
\section{Reformulierung des Driftarguments}
\label{sec:drift_korrektur}
% ====================================================

\subsection{Korrekte Erwartungswerte}

Die numerischen Ergebnisse aus Abschnitt~\ref{sec:numerik} legen nahe,
dass die $\nu$-Verteilung konditioniert auf $l = k$ \emph{dieselbe} ist
wie ohne Konditionierung auf $l$:
\begin{align*}
  P(\nu = j \mid l = 0) &\approx 2^{-(j-1)}, \quad j \geq 2, \\
  P(\nu = j \mid l = k) &\approx 2^{-(j-1)}, \quad j \geq 2, \quad k \geq 1.
\end{align*}
In beiden Fällen gilt $\E[\nu \mid l{=}k] = 3$.

\begin{bemerkung}
Dies ist heuristisch plausibel: Wenn $m$ (modulo Potenzen von $2$) equidistributed
ist und die Bedingung $l = k$ keine signifikante Information über $m \pmod{2^r}$
für großes $r$ liefert, dann verhält sich $\vii(3^{k+1}m - 1)$ wie eine
geometrische Variable auf $\{1, 2, 3, \ldots\}$, verschoben um $1$.
\end{bemerkung}

\subsection{Korrigierte Driftformel}

Mit $\E[\nu \mid l{=}k] = \alpha$ für $k \geq 0$ (hier: $\alpha \approx 3$):
\[
  \E[\logF \mid l = k]
  = (k+1)\log_2 3 - (k + \alpha)
  = k(\log_2 3 - 1) + (\log_2 3 - \alpha).
\]
Mit $\log_2 3 \approx 1{,}585$, also $\log_2 3 - 1 \approx 0{,}585$
und $\log_2 3 - \alpha \approx 1{,}585 - 3 = -1{,}415$:
\[
  \E[\logF \mid l = k] \approx 0{,}585\,k - 1{,}415.
\]
Dies ist identisch mit dem $l = 0$-Wert ($k = 0$: $-1{,}415$)
und wächst mit $k$ --- aber die Gewichtung mit $P(l{=}k) = 2^{-(k+1)}$ dominiert.

\begin{proposition}[Korrigierter Gesamtdrift]
\label{prop:korrekter_drift}
Mit $\E[\nu \mid l{=}k] \approx 3$ für alle $k \geq 0$:
\begin{align*}
  \E[\logF]
  &= \sum_{k=0}^{\infty} P(l{=}k) \cdot \E[\logF \mid l{=}k] \\
  &\approx \sum_{k=0}^{\infty} \frac{1}{2^{k+1}} \cdot (0{,}585\,k - 1{,}415) \\
  &= 0{,}585 \underbrace{\sum_{k=0}^{\infty} k \cdot 2^{-(k+1)}}_{= \,1}
    - 1{,}415 \underbrace{\sum_{k=0}^{\infty} 2^{-(k+1)}}_{= \,1}
  = 0{,}585 - 1{,}415 = -0{,}830.
\end{align*}
\end{proposition}

\begin{tcolorbox}[hauptsatz, title={Korrigiertes Hauptresultat (heuristisch)}]
Mit der numerisch gestützten Annahme $\E[\nu \mid l{=}k] \approx 3$
für \emph{alle} $k \geq 0$ ergibt die konditionierte Driftrechnung:
\[
  \boxed{\E[\logF] \approx -0{,}830 < 0.}
\]
Dies ist \emph{robuster} und \emph{negativer} als der fehlerhafte Wert
$-0{,}33$ aus dem Ausgangsdokument.
Insbesondere ist das Ergebnis unabhängig von der Kettenlänge $k$,
da $\E[\nu \mid l{=}k]$ offenbar nicht von $k$ abhängt.
\end{tcolorbox}

\subsection{Korrigierte Sensitivitätsanalyse}

\begin{table}[h!]
\centering
\caption{Korrigierte Sensitivitätsanalyse: $\E[\logF]$ als Funktion von $\alpha = \E[\nu]$}
\label{tab:sensitivitaet_korr}
\renewcommand{\arraystretch}{1.3}
\begin{tabular}{
  S[table-format=1.1]
  S[table-format=+1.4]
  l
}
\toprule
{$\alpha$} & {$\E[\logF](\alpha)$} & Interpretation \\
\midrule
2{,}0 & -0{,}415 & \textit{(Dokument-Fehlannahme: E[nu]=2)} \\
2{,}5 & -0{,}665 & schwach kontraktiv \\
3{,}0 & -0{,}830 & \textbf{korrekt (numerische Schätzung)} \\
3{,}5 & -1{,}080 & stärker kontraktiv \\
4{,}0 & -1{,}330 & sehr stark kontraktiv \\
\midrule
{1{,}585} & {0{,}000} & \textbf{kritischer Wert} $\alpha^* = \log_2 3$ \\
\bottomrule
\end{tabular}
\end{table}

\begin{bemerkung}
Mit der korrekten Formel für die uniformen $P(l{=}k)$ lautet der kritische Wert:
\[
  \alpha^* = \log_2 3 \approx 1{,}585,
\]
da $\E[\logF](\alpha) = 0$ iff $\alpha = \log_2 3$.
Für $\alpha = \E[\nu \mid l{=}k] \approx 3 > 1{,}585 = \alpha^*$ gilt
also $\E[\logF] < 0$ mit erheblichem Abstand.

Übrigens: Die im Ausgangsdokument genannte Formel $\E[\logF](\alpha) = 0{,}670 - \alpha/2$
ist wegen der falschen Zerlegung von $l{=}0$ vs.\ $l{\geq}1$ auch formal nicht korrekt;
die richtige allgemeine Formel lautet $\E[\logF](\alpha) = \log_2 3 - \alpha$.
\end{bemerkung}

% ====================================================
\section{Was noch bewiesen werden müsste}
\label{sec:offene_probleme}
% ====================================================

\begin{tcolorbox}[lücke, title={Offenes Problem 1: Maßstruktur und Equidistribution}]
\textbf{Was fehlt:}
Ein Beweis, dass die $m$-Werte $(n_0 + 1)/2^{l+1}$ aus C-Ketten der Länge $k$
equidistributed modulo $2^r$ sind (für alle $r$).

\medskip
\textbf{Konkret:} Zeige, dass für alle $r \geq 1$ und alle ungeraden $a$ mit
$\gcd(a, 2) = 1$:
\[
  \lim_{X \to \infty}
  \frac{\#\{n_0 \leq X : n_0 \equiv 11 \pmod{12},\, l(n_0)=k,\, m(n_0) \equiv a \pmod{2^r}\}}
  {\#\{n_0 \leq X : n_0 \equiv 11 \pmod{12},\, l(n_0)=k\}}
  = 2^{-(r-1)}.
\]

\textbf{Konsequenz:} Equidistribution würde $P(\nu{=}j \mid l{=}k) = 2^{-(j-1)}$ für
$j \geq 2$ beweisen, und damit $\E[\nu \mid l{=}k] = 3$ rigoros begründen.
\end{tcolorbox}

\begin{tcolorbox}[lücke, title={Offenes Problem 2: Stationäre Bahnverteilung}]
\textbf{Was fehlt:}
Nachweis, dass die Block-Längen $l(B_1), l(B_2), \ldots$ entlang einer
Collatz-Bahn die geometrische Verteilung $P(l{=}k) = 2^{-(k+1)}$
(im Zeitmittel) annehmen.

\medskip
\textbf{Konkret:} Für fast alle $n \in \N$:
\[
  \lim_{N \to \infty} \frac{1}{N} \sum_{i=1}^{N} \mathbf{1}[l(B_i) = k]
  \stackrel{?}{=} 2^{-(k+1)}.
\]

\textbf{Konsequenz:} Dies würde zusammen mit Op.\ 1 die heuristische
Driftrechnung zu einem (fast sicheren) Aussage über Zeitmittel erheben.
Ohne Ergodizitätssatz für die Collatz-Abbildung bleibt dies offen.
\end{tcolorbox}

\begin{tcolorbox}[lücke, title={Offenes Problem 3: Gleichmäßige Negativität}]
\textbf{Was fehlt:}
Selbst wenn $\E[\logF] \approx -0{,}83 < 0$ (im Ensemblemittel oder Zeitmittel)
bewiesen wäre, bliebe für einen vollständigen Collatz-Beweis nötig:
\begin{enumerate}[noitemsep]
  \item \emph{Ausschluss nicht-trivialer Zyklen:}
        Ein Zyklus $n_1 \to n_2 \to \cdots \to n_r \to n_1$ hätte
        $\sum_i \logF(B_i) = 0$, was einem lokal positiven Drift nicht
        widerspricht.
  \item \emph{Ausschluss divergenter Bahnen:}
        Negativer mittlerer Drift schließt $T^k(n) \to \infty$ nicht
        für alle $n$ aus, sondern nur für ``generische'' $n$.
\end{enumerate}
\end{tcolorbox}

\begin{table}[h!]
\centering
\caption{Korrigierte Statusübersicht der Argumentationsschritte}
\label{tab:status_korr}
\renewcommand{\arraystretch}{1.4}
\begin{tabular}{p{5.5cm} c p{6.5cm}}
\toprule
\textbf{Aussage} & \textbf{Status} & \textbf{Anmerkung} \\
\midrule
$\nu \geq 2$ für alle $l \geq 1$
  & \textcolor{bewiesengrün}{\textbf{bewiesen}}
  & Triviales Paritätsargument (Satz~\ref{satz:nu_geq_2}) \\
$3 \mid m$ für $n_0 \equiv 11 \pmod{12}$
  & \textcolor{bewiesengrün}{\textbf{bewiesen}}
  & Direkte Rechnung (Satz~\ref{satz:drei_teilt_m}) \\
$s$-Verteilung geometrisch
  & \textcolor{warnorange}{\textbf{Konjektur}}
  & Numerisch exakt; analytisch offen \\
$\E[\nu \mid l{=}k] \approx 3$
  & \textcolor{warnorange}{\textbf{heuristisch}}
  & Numerisch für $n_0 \leq 10\,000$ bestätigt \\
$P(l{=}k) = 2^{-(k+1)}$ auf Bahnen
  & \textcolor{warnorange}{\textbf{heuristisch}}
  & Gilt für uniform gewählte $n_0$, nicht auf Bahnen \\
$\E[\logF] \approx -0{,}83$
  & \textcolor{warnorange}{\textbf{heuristisch}}
  & Folgt aus obigen Annahmen \\
Zeitmittel $\frac{1}{N}\sum \logF \to -0{,}83$
  & \textcolor{lückenrot}{\textbf{offen}}
  & Erfordert Ergodizität \\
Zyklenfreiheit / Nicht-Divergenz
  & \textcolor{lückenrot}{\textbf{offen}}
  & Unabhängig vom Driftargument \\
\bottomrule
\end{tabular}
\end{table}

% ====================================================
\appendix
\section{Verifikation: LTE-Formel für \texorpdfstring{$\vii(3^N - 1)$}{nu2(3^N-1)}}
\label{app:lte}
% ====================================================

Zum Vergleich mit dem Ausgangsdokument: Die exakte LTE-Formel für
$\vii(3^N - 1)$ (Spezialfall $m' = 1$):
\begin{align*}
  N \text{ ungerade:} &\quad \vii(3^N - 1) = 1, \\
  N \text{ gerade:}   &\quad \vii(3^N - 1) = 2 + \vii(N).
\end{align*}
Dies folgt aus dem LTE für $p=2$: Da $2 \mid (3-1)$, gilt für ungerades $N$:
$\vii(3^N - 1) = \vii(3-1) = 1$.
Für gerades $N$: $\vii(3^N - 1) = \vii(3-1) + \vii(3+1) + \vii(N) - 1 = 1 + 2 + \vii(N) - 1 = 2 + \vii(N)$.

\medskip
Für allgemeines $m' \neq 1$ ist $3^N m' - 1$ \emph{nicht} von der Form $a^N - b^N$
(mit konstantem $a$), weshalb das LTE nicht direkt anwendbar ist.
Die Schranke $\vii(3^{l+1} m - 1) \leq \vii(3m-1) + \vii(l+1)$ im Ausgangsdokument
ist daher als heuristische Übertragung zu werten, nicht als strikte Anwendung des LTE.

% ====================================================
\section{Verbindung zu \textit{collatz\_lyapunov\_konditioniert.tex}}
\label{app:verbindung}
% ====================================================

Dieses Dokument korrigiert und erweitert
\textit{collatz\_lyapunov\_konditioniert.tex}:

\begin{itemize}[noitemsep]
  \item \textbf{Korrektur:} $\E[\nu \mid l{=}k] \approx 3$ (nicht $2$);
        Driftformel: $\E[\logF] \approx -0{,}83$ (nicht $-0{,}33$).
  \item \textbf{Neues Ergebnis:} $3 \mid m$ für $n_0 \equiv 11 \pmod{12}$
        (Satz~\ref{satz:drei_teilt_m}).
  \item \textbf{Neue Konjektur:} $s = \viii(m)$ ist geometrisch mit
        $P(s{=}j) = \frac{2}{3}(1/3)^{j-1}$.
  \item \textbf{Klarstellung:} $\nu \geq 2$ ist trivial (Parität),
        kein LTE-Ergebnis.
  \item \textbf{Forschungsfrage:} Equidistribution der $m$-Werte modulo
        $2^r$ --- dies ist der entscheidende fehlende Schritt für eine
        rigorose Driftrechnung.
\end{itemize}

Die grundsätzliche Botschaft des Ausgangsdokuments bleibt erhalten:
Der heuristische Drift ist negativ, $\E[\logF] < 0$.
Er ist sogar \emph{deutlich} negativer als ursprünglich behauptet.
Die fundamentalen Lücken (Ergodizität, Zyklenfreiheit) sind dieselben.

\end{document}
